\documentclass[11pt]{article}
\usepackage[a4paper,margin=2.5cm]{geometry}
\usepackage{amsmath,amssymb,mathtools}
\usepackage[T1]{fontenc}
\usepackage[utf8]{inputenc}
\usepackage{lmodern}
\usepackage{microtype}
\usepackage{xcolor}
\definecolor{grey}{rgb}{0.8,0.8,0.8}

\newcommand{\dd}{\mathrm d}
\newcommand{\ii}{\mathrm i}
\newcommand{\eps}{\epsilon}
\newcommand{\Li}{\operatorname{Li}}
\newcommand{\sgn}{\operatorname{sgn}}

\begin{document}

\section{Two-loop Differential Collinear Kernel}


We consider the double collinear building block.  The two masses inside each collinear pair are kept distinct:
\begin{equation}
\mathcal I(x,y)=
\int \frac{d^d k}{i\pi^{d/2}}\frac{d^d l}{i\pi^{d/2}}\,
\frac{\delta\!\left(x-l\!\cdot\!\eta\right)\delta\!\left(y-k\!\cdot\!\eta\right)}
{[l^2-m_l^2+i0][(l-p)^2-M_l^2+i0][k^2-m_k^2+i0][(l+k-p)^2-M_k^2+i0]}.
\label{eq:fullI}
\end{equation}
We normalize
\begin{equation}
p\!\cdot\!\eta=1,
\qquad \eta^2=0.
\label{eq:lightcone-normalization}
\end{equation}
Define
\begin{equation}
K(p-l;y)\equiv
\int \frac{d^d k}{i\pi^{d/2}}\,
\frac{\delta\!\left(y-k\!\cdot\!\eta\right)}
{[k^2-m_k^2+i0][(k-p+l)^2-M_k^2+i0]}.
\label{eq:Kdef}
\end{equation}
Then
\begin{equation}
\mathcal I(x,y)=
\int \frac{d^d l}{i\pi^{d/2}}\,
\frac{\delta\!\left(x-l\!\cdot\!\eta\right)}{[l^2-m_l^2+i0][(l-p)^2-M_l^2+i0]}
\,K(p-l;y).
\label{eq:I-via-K}
\end{equation}



\subsubsection{Light-cone decomposition}


\begin{equation}
l=xp+\beta_l\eta+l_\perp,
\qquad
k=yp+\beta_k\eta+k_\perp,
\label{eq:lck}
\end{equation}
with
\begin{equation}
p\cdot l_\perp=\eta\cdot l_\perp=0,
\qquad
p\cdot k_\perp=\eta\cdot k_\perp=0.
\label{eq:transverse-orthogonality}
\end{equation}
Then indeed
\begin{equation}
l\cdot\eta=x,
\qquad
k\cdot\eta=y.
\label{eq:delta-projections}
\end{equation}
In these variables,
\begin{equation}
d^d l=dx\,d\beta_l\,d^{d-2}l_\perp,
\qquad
d^d k=dy\,d\beta_k\,d^{d-2}k_\perp.
\label{eq:lightcone-measures}
\end{equation}
\subsubsection{The two $k$-dependent denominators}

The first denominator is
\begin{equation}
k^2=y^2p^2+2y\beta_k+k_\perp^2,
\qquad
D_1\equiv k^2-m_k^2+i0=2y\beta_k+\bigl(y^2p^2+k_\perp^2-m_k^2\bigr)+i0.
\label{eq:D1}
\end{equation}
For the second denominator we use
\[
k-p+l=(x+y-1)p+(\beta_k+\beta_l)\eta+(k_\perp+l_\perp),
\]
so that
\begin{equation}
(k-p+l)^2=(x+y-1)^2p^2+2(x+y-1)(\beta_k+\beta_l)+(k_\perp+l_\perp)^2,
\label{eq:kplsq}
\end{equation}
and hence
\begin{equation}
D_2\equiv (k-p+l)^2-M_k^2+i0
=2(x+y-1)(\beta_k+\beta_l)+\bigl((x+y-1)^2p^2+(k_\perp+l_\perp)^2-M_k^2\bigr)+i0.
\label{eq:D2}
\end{equation}
Equivalently, with
\[
q=p-l=(1-x)p-\beta_l\eta-l_\perp,
\qquad q\cdot\eta=1-x,
\]
one may write
\begin{equation}
D_2
=2\bigl(y-(1-x)\bigr)(\beta_k+\beta_l)
+\bigl((y-(1-x))^2p^2+(k_\perp+l_\perp)^2-M_k^2\bigr)+i0.
\label{eq:D2-alt}
\end{equation}
Thus
\begin{equation}
K(p-l;y)=
\int \frac{d\beta_k\,d^{d-2}k_\perp}{i\pi^{d/2}}\,
\frac{1}{D_1D_2}.
\label{eq:Kbetak}
\end{equation}

\subsubsection{Pole structure in $\beta_k$ and support in $y$}

The two poles in the complex $\beta_k$ plane are
\begin{equation}
\beta_k^{(1)}=
\frac{m_k^2-k_\perp^2-y^2p^2-i0}{2y},
\label{eq:beta1}
\end{equation}
and
\begin{equation}
\beta_k^{(2)}=
-\beta_l+
\frac{M_k^2-(k_\perp+l_\perp)^2-(y-(1-x))^2p^2-i0}{2(y-(1-x))}.
\label{eq:beta2}
\end{equation}
The $\beta_k$ integral is nonzero only when the two poles lie on opposite sides of the real axis. Assuming $1-x>0$, this happens precisely for
\begin{equation}
0<y<1-x.
\label{eq:support-y}
\end{equation}
Therefore the subloop carries the support
\begin{equation}
\Theta(0\le y\le 1-x).
\label{eq:support-theta}
\end{equation}
% This is the exact analogue of the support $\Theta(0\le x\le 1)$ in Eq.~(3.20) of Anastasiou--Sterman.

\subsubsection{The $\beta_k$ integral by residues}

We evaluate the $\beta_k$ integral by taking the residue at the pole of $D_1$. Since the pole is simple,
\[
\operatorname{Res}_{\beta_k=\beta_k^{(1)}}\frac{1}{D_1D_2}
=
\frac{1}{D_1'(\beta_k^{(1)})\,D_2(\beta_k^{(1)})},
\qquad D_1'(\beta_k)=2y.
\]
Hence the residue contributes the factor $1/(2y)$.
\\
At the pole $\beta_k=\beta_k^{(1)}$ one finds
\begin{align}
D_2\bigl(\beta_k^{(1)}\bigr)
={}&
\frac{y-(1-x)}{y}\bigl(m_k^2-k_\perp^2-y^2p^2\bigr)
+2\bigl(y-(1-x)\bigr)\beta_l
\nonumber\\
&
+(y-(1-x))^2p^2+(k_\perp+l_\perp)^2-M_k^2+i0.
\label{eq:D2atpole1}
\end{align}
Using $(k_\perp+l_\perp)^2=k_\perp^2+2k_\perp\cdot l_\perp+l_\perp^2$, this becomes
\begin{align}
D_2\bigl(\beta_k^{(1)}\bigr)
={}&
-\frac{1-x}{y}k_\perp^2+2k_\perp\cdot l_\perp
\nonumber\\
&+
\left[
\frac{y-(1-x)}{y}m_k^2+2\bigl(y-(1-x)\bigr)\beta_l+l_\perp^2-(1-x)(x+y-1)p^2-M_k^2
\right]
+i0.
\label{eq:D2atpole2}
\end{align}
Next we use the exact identity
\begin{equation}
(p-l)^2=(1-x)^2p^2-2(1-x)\beta_l+l_\perp^2,
\label{eq:plsq}
\end{equation}
which implies
\[
l_\perp^2=(p-l)^2-(1-x)^2p^2+2(1-x)\beta_l.
\]
After a short algebraic rearrangement, it is convenient to introduce
\begin{equation}
\lambda\equiv \frac{y}{1-x},
\qquad
\mu_k^2(\lambda)\equiv (1-\lambda)m_k^2+\lambda M_k^2.
\label{eq:lambda}
\end{equation}
and one obtains
\begin{equation}
D_2\bigl(\beta_k^{(1)}\bigr)
=-\frac{1-x}{y}
\left[
\bigl(k_\perp-\lambda l_\perp\bigr)^2+\Delta_\perp-i0
\right],
\label{eq:D2factorized}
\end{equation}
where
\begin{equation}
\Delta_\perp
=
\mu_k^2(\lambda)-\lambda(1-\lambda)(p-l)^2-\lambda\bigl(1-\lambda x\bigr)p^2,
\label{eq:Deltaperp}
\end{equation}
% The $-i0$ inside the square bracket is the correct prescription after factoring out the negative real number $-(1-x)/y$.
Therefore
\begin{equation}
\frac{1}{D_2\bigl(\beta_k^{(1)}\bigr)}
=-\frac{y}{1-x}
\frac{1}{\bigl(k_\perp-\lambda l_\perp\bigr)^2+\Delta_\perp-i0}.
\label{eq:invD2}
\end{equation}
Combining this with the residue factor $1/(2y)$, the factors of $y$ cancel and one gets
\begin{equation}
K(p-l;y)=
\frac{\Theta(0\le y\le 1-x)}{1-x}
\int \frac{d^{d-2}k_\perp}{\pi^{(d-2)/2}}
\frac{1}{\bigl(k_\perp-\lambda l_\perp\bigr)^2+\Delta_\perp-i0}.
\label{eq:K-after-residue}
\end{equation}
Now shift
\[
r_\perp=k_\perp-\lambda l_\perp,
\qquad d^{d-2}k_\perp=d^{d-2}r_\perp,
\]
and use the standard Euclidean integral in $n=d-2=2-2\epsilon$ dimensions,
\[
\int \frac{d^n r}{\pi^{n/2}}\frac{1}{(r^2+\Delta)^\alpha}
=
\frac{\Gamma(\alpha-n/2)}{\Gamma(\alpha)}\,\Delta^{n/2-\alpha}.
\]
With $\alpha=1$ this gives
\begin{equation}
K(p-l;y)=
\Gamma(\epsilon)
\frac{\Theta(0\le y\le 1-x)}{1-x}
\Bigl[\mu_k^2(\lambda)-\lambda(1-\lambda)(p-l)^2-\lambda(1-\lambda x)p^2-i0\Bigr]^{-\epsilon}.
\label{eq:K-final}
\end{equation}
Equivalently,
\begin{equation}
K(p-l;y)
=
\Gamma(\eps)
\frac{\Theta(0\le y\le 1-x)}{1-x}
\left[C(l)-\ii0\right]^{-\eps},
\label{eq:k-subloop}
\end{equation}
where
\begin{equation}
C(l)
=
\mu_k^2(\lambda)-\lambda(1-\lambda)(p-l)^2
-\lambda(1-\lambda x)p^2.
\label{eq:C-l}
\end{equation}
Thus the full integral is reduced to
\begin{equation}
\boxed{
\mathcal I(x,y)
=
\Gamma(\eps)
\frac{\Theta(0\le y\le 1-x)}{1-x}
\int \frac{\dd^d l}{\ii\pi^{d/2}}
\frac{\delta(x-l\cdot\eta)}
{[l^2-m_l^2+\ii0][(l-p)^2-M_l^2+\ii0]}
[C(l)-\ii0]^{-\eps}.
}
\label{eq:after-k-subloop}
\end{equation}
% \begin{equation}
% K(p-l;y)=
% \Gamma(\epsilon)
% \frac{\Theta(0\le y\le 1-x)}{1-x}
% \left[
% \right]^{-\epsilon}.
% \end{equation}

% Thus
% \begin{align}
% \mathcal I(x,y)
% ={}&
% \Gamma(\epsilon)\,
% \frac{\Theta(0\leq y\leq 1-x)}{1-x}
% \int \frac{\dd^d l}{\ii\pi^{d/2}}
% \frac{\delta(x-l\!\cdot\!\eta)}{[l^2-m_l^2+\ii0][(l-p)^2-M_l^2+\ii0]}
% \nonumber\\
% &\times
% \left[
% \right]^{-\epsilon}.
% \end{align}
\newpage
\subsection{Method 1: Simultaneous Parametrization}

\subsubsection{Combine all $l$-dependent factors simultaneously}
Introduce
\begin{equation}
A=l^2-m_l^2+i0,\qquad B=(l-p)^2-M_l^2+i0.
\label{eq:method1-A-B-def}
\end{equation}
Then
\begin{equation}
\frac{\Gamma(\epsilon)}{A\,B\,C^{\epsilon}}=
\Gamma(2+\epsilon)\int du\,dv\,dw\,\delta(1-u-v-w)
\frac{w^{\epsilon-1}}{[uA+vB+wC]^{2+\epsilon}}.
\label{eq:method1-feynman-param}
\end{equation}
Hence
\begin{equation}
\mathcal I(x,y)=
\frac{\Gamma(2+\epsilon)\Theta(0\le y\le 1-x)}{1-x}
\int_{\Delta_2}du\,dv\,dw\,w^{\epsilon-1}
\int\frac{d^d l}{i\pi^{d/2}}
\frac{\delta(x-l\cdot\eta)}{D(l)^{2+\epsilon}},
\label{eq:method1-parametric-before-l}
\end{equation}
where $u,v,w\ge0$, $u+v+w=1$ and
\begin{equation}
D(l)=uA+vB+wC.
\label{eq:method1-D-def}
\end{equation}

\subsubsection{Quadratic form in $l$}
Writing
\begin{equation}
a=\lambda(1-\lambda),\qquad b=\lambda(1-\lambda x),
\label{eq:method1-a-b-def}
\end{equation}
and
\begin{equation}
\mu_l^2(x)\equiv (1-x)m_l^2+xM_l^2,
\qquad
\mu_l^2(u)\equiv u m_l^2+(1-u)M_l^2,
\label{eq:mu-l-def}
\end{equation}
one finds after expansion
\begin{equation}
D(l)=U l^2-2V\,l\cdot p+W+i0,
\label{eq:method1-quadratic-form}
\end{equation}
with
\begin{align}
U&=u+v-aw,\\
V&=v-aw,\\
W&=[v-w(a+b)]p^2-u m_l^2-v M_l^2+w\mu_k^2(\lambda).
\label{eq:method1-UVW-def}
\end{align}
Shift
\[
l=q+\frac{V}{U}p,
\]
so that
\begin{equation}
D(l)=U q^2+\Delta(u,v,w)+i0,
\label{eq:method1-shifted-denominator}
\end{equation}
where
\begin{equation}
\Delta=W-\frac{V^2}{U}p^2.
\label{eq:method1-Delta-def}
\end{equation}

\subsubsection{Delta function by Fourier representation}
Use
\[
\delta(x-l\cdot\eta)=\int\frac{d\omega}{2\pi}e^{i\omega(x-l\cdot\eta)}.
\]
Since $p\cdot\eta=1$ and $\eta^2=0$, the Gaussian $q$-integration gives
\begin{equation}
\int \frac{d^d q}{i\pi^{d/2}}\frac{e^{-i\omega q\cdot\eta}}{(Uq^2+\Delta)^{2+\epsilon}}
=U^{-d/2}\frac{\Gamma(2\epsilon)}{\Gamma(2+\epsilon)}\Delta^{-2\epsilon},
\end{equation}
independent of $\omega$. The Fourier source in the shifted \(q\)-integral does not generate an
\(\omega\)-dependent factor because the source vector is lightlike. To see this,
use a Schwinger representation,
\[
\frac{1}{(Uq^2+\Delta)^{2+\epsilon}}
=
\frac{1}{\Gamma(2+\epsilon)}
\int_0^\infty d\alpha\,
\alpha^{1+\epsilon}
e^{-\alpha(Uq^2+\Delta)} .
\]
Then the \(q\)-dependent part is
\[
\int d^d q\,
\exp\!\left[-\alpha U q^2-i\omega q\cdot\eta\right].
\]
Completing the square gives the standard Gaussian result
\[
\int d^d q\,
\exp\!\left[-\alpha U q^2-i\omega q\cdot\eta\right]
=
\left(\frac{\pi}{\alpha U}\right)^{d/2}
\exp\!\left[-\frac{\omega^2\eta^2}{4\alpha U}\right].
\]
Since the auxiliary vector is lightlike,
\[
\eta^2=0,
\]
the source-dependent exponential becomes
\[
\exp\!\left[-\frac{\omega^2\eta^2}{4\alpha U}\right]=1.
\]
Therefore the Fourier source does not modify the Gaussian integral:
\[
\int d^d q\,
\exp\!\left[-\alpha U q^2-i\omega q\cdot\eta\right]
=
\left(\frac{\pi}{\alpha U}\right)^{d/2}.
\]
Consequently,
\[
\int \frac{d^d q}{i\pi^{d/2}}\,
\frac{e^{-i\omega q\cdot\eta}}
{(Uq^2+\Delta)^{2+\epsilon}}
=
\int \frac{d^d q}{i\pi^{d/2}}\,
\frac{1}
{(Uq^2+\Delta)^{2+\epsilon}},
\]
Hence
\[
\int \frac{d^d q}{i\pi^{d/2}}\,
\frac{e^{-i\omega q\cdot\eta}}
{(Uq^2+\Delta)^{2+\epsilon}}
=
U^{-d/2}
\frac{\Gamma(2\epsilon)}{\Gamma(2+\epsilon)}
\Delta^{-2\epsilon}.
\]
The essential point is that the source is proportional to \(\eta^\mu\), and
therefore its square is proportional to \(\eta^2\). Since \(\eta^2=0\), the
Fourier source leaves no \(\omega\)-dependence behind.Therefore the $\omega$ integral returns
\[
\delta\!\left(x-\frac{V}{U}\right).
\]
Thus
\begin{equation}
\mathcal I(x,y)=
\frac{\Gamma(2\epsilon)\Theta(0\le y\le 1-x)}{1-x}
\int_{\Delta_2}du\,dv\,dw\,w^{\epsilon-1}U^{-d/2}
\delta\!\left(x-\frac{V}{U}\right)
\Delta^{-2\epsilon}.
\end{equation}

\subsubsection{Collapse to one parameter}
Using $u+v+w=1$, one has
\[
U=1-(1+a)w,
\]
which is independent of $v$. The delta functions fix
\[
v=aw+xU,\qquad u=(1-x)U.
\]
Hence
\begin{equation}
\mathcal I(x,y)=
\frac{\Gamma(2\epsilon)\Theta(0\le x\le1)\Theta(0\le y\le1-x)}{1-x}
\int_0^{1/(1+a)}dw\,w^{\epsilon-1}U^{-1+\epsilon}[\Delta_x(w)-i0]^{-2\epsilon},
\label{eq:method1-w-integral-before-euler}
\end{equation}
with
\[
U=1-(1+a)w.
\]
The linear function $\Delta_x(w)$ can be written as
\begin{equation}
\Delta_x(w)=\Delta_0+w\Delta_w,
\label{eq:method1-Delta-x-def}
\end{equation}
where
\begin{align}
\Delta_0 &= \mu_l^2(x)-x(1-x)p^2,\\
\Delta_w &= \bigl((1+a)x(1-x)+b\bigr)p^2
-\mu_k^2(\lambda)+aM_l^2-(1+a)\mu_l^2(x).
\label{eq:method1-Delta0-Deltaw-def}
\end{align}

\subsubsection{Map to Euler form}
Set
\begin{equation}
t=(1+a)w,
\label{eq:method1-euler-map}
\end{equation}
so $t\in[0,1]$. Then
\begin{equation}
\mathcal I(x,y)=
\frac{\Gamma(2\epsilon)\Theta}{1-x}(1+a)^{-\epsilon}(\Delta_0-i0)^{-2\epsilon}
\int_0^1dt\,t^{\epsilon-1}(1-t)^{\epsilon-1}(1-zt)^{-2\epsilon},
\end{equation}
where
\begin{equation}
\Theta\equiv \Theta(0\le x\le1)\Theta(0\le y\le1-x),
\qquad
z=-\frac{\Delta_w}{(1+a)\Delta_0}.
\label{eq:method1-theta-z-def}
\end{equation}
Using Euler's integral,
\[
\int_0^1dt\,t^{\alpha-1}(1-t)^{\beta-1}(1-zt)^{-\rho}
=B(\alpha,\beta)\,{}_2F_1(\rho,\alpha;\alpha+\beta;z),
\]
with $\alpha=\beta=\epsilon$, $\rho=2\epsilon$, gives
\begin{equation}
\mathcal I(x,y)=
\frac{\Gamma(\epsilon)^2}{1-x}\Theta
(1+a)^{-\epsilon}(\Delta_0-i0)^{-2\epsilon}
{}_2F_1(2\epsilon,\epsilon;2\epsilon;z).
\label{eq:method1-hypergeometric-form}
\end{equation}
Now use
\[
{}_2F_1(A,B;A;z)=(1-z)^{-B},
\]
so that
\[
{}_2F_1(2\epsilon,\epsilon;2\epsilon;z)=(1-z)^{-\epsilon}.
\]
This combines with the prefactors into a second scale $\Delta_1$.  With the present sign convention for $\Delta_x(w)$,
\begin{equation}
\Delta_1=(1+a)\Delta_0+\Delta_w.
\label{eq:method1-Delta1-from-Deltaw}
\end{equation}

\subsubsection{Final exact kernel}
We obtain
\begin{equation}
\boxed{
\mathcal I(x,y)=
\frac{\Gamma(\epsilon)^2}{1-x}
\Theta(0\le x\le1)\Theta(0\le y\le1-x)
[\Delta_0-i0]^{-\epsilon}[\Delta_1-i0]^{-\epsilon}
}
\label{eq:closed-form-kernel}
\end{equation}
with
\begin{equation}
\Delta_0=\mu_l^2(x)-x(1-x)p^2,
\label{eq:Delta0-final}
\end{equation}
and
\begin{equation}
\Delta_1=
-\mu_k^2(\lambda)+\lambda(1-\lambda)M_l^2+\lambda(1-\lambda x)p^2,
\qquad \lambda=\frac{y}{1-x}.
\label{eq:Delta1-final}
\end{equation}
Equivalently,
\begin{equation}
\Delta_1=-
\left(1-\frac{y}{1-x}\right)m_k^2
-\frac{y}{1-x}M_k^2
+\frac{y}{1-x}\left(1-\frac{y}{1-x}\right)M_l^2
+\frac{y}{1-x}\left(1-\frac{xy}{1-x}\right)p^2.
\label{eq:Delta1-final-expanded}
\end{equation}



\newpage

\subsection{Method 2: Consecutive Parametrization and Light-Cone Decomposition}

\subsubsection{Combining the two quadratic $l$ propagators}

Define
\begin{equation}
A=l^2-m_l^2+\ii0,
\qquad
B=(l-p)^2-M_l^2+\ii0.
\label{eq:method2-A-B-def}
\end{equation}
We combine them by
\begin{equation}
\frac{1}{AB}=\int_0^1 \dd u\,
\frac{1}{[uA+(1-u)B]^2}.
\label{eq:method2-feynman-param}
\end{equation}
A short calculation gives
\begin{equation}
uA+(1-u)B=\bigl(l-(1-u)p\bigr)^2-\mu_l^2(u)+u(1-u)p^2+\ii0.
\label{eq:method2-combined-denominator}
\end{equation}
We therefore shift
\begin{equation}
L\equiv l-(1-u)p,
\qquad
l=L+(1-u)p,
\qquad
p-l=up-L.
\label{eq:method2-L-shift}
\end{equation}
The delta function becomes
\begin{equation}
\delta(x-l\!\cdot\!\eta)=
\delta\bigl(L\!\cdot\!\eta-(x-1+u)\bigr).
\label{eq:method2-delta-shift}
\end{equation}
Hence
\begin{align}
\mathcal I(x,y)
={}&
\Gamma(\epsilon)\,
\frac{\Theta(0\leq y\leq 1-x)}{1-x}
\int_0^1 \dd u
\int \frac{\dd^d L}{\ii\pi^{d/2}}
\delta\bigl(L\!\cdot\!\eta-(x-1+u)\bigr)
\nonumber\\
&\times
\frac{[\mathcal M_u(L;x,y)-\ii0]^{-\epsilon}}
{[L^2-\mu_l^2(u)+u(1-u)p^2+\ii0]^2},
\label{eq:Ishifted}
\end{align}
where
\begin{equation}
\mathcal M_u(L;x,y)=
\mu_k^2(\lambda)-\lambda(1-\lambda)(L^2-2uL\!\cdot p+u^2p^2)-\lambda(1-\lambda x)p^2.
\label{eq:MuL}
\end{equation}

\subsubsection{The remaining light-cone decomposition}

Now decompose
\begin{equation}
L=\xi p+\alpha\eta+L_\perp,
\qquad
L\!\cdot\!\eta=\xi,
\qquad
L\!\cdot p=\alpha+\xi p^2,
\qquad
L^2=2\xi\alpha+\xi^2p^2+L_\perp^2.
\label{eq:method2-L-lightcone}
\end{equation}
The delta function fixes
\begin{equation}
\xi=x-1+u.
\label{eq:method2-xi-fixed}
\end{equation}
We introduce
\begin{equation}
a_u\equiv u-1+x.
\label{eq:audef}
\end{equation}
Define also
\begin{equation}
X\equiv 1-x,
\qquad
\kappa\equiv\lambda(1-\lambda)=\frac{y(1-x-y)}{(1-x)^2}.
\label{eq:X-kappa-def}
\end{equation}
Then the squared denominator is affine in $\alpha$:
\begin{equation}
L^2-\mu_l^2(u)+u(1-u)p^2+\ii0
=2a_u\alpha+\Lambda_u(L_\perp;x)+\ii0,
\label{eq:quadalpha}
\end{equation}
with
\begin{equation}
\Lambda_u(L_\perp;x)=a_u^2p^2+L_\perp^2-\mu_l^2(u)+u(1-u)p^2.
\label{eq:Lambdau}
\end{equation}
The noninteger factor is also affine in $\alpha$:
\begin{equation}
\mathcal M_u(\alpha,L_\perp;x,y)=\beta_u(x,y)\alpha+\rho_u(L_\perp;u,x,y),
\label{eq:Muaffine}
\end{equation}
where
\begin{equation}
\beta_u(x,y)=2\lambda(1-\lambda)(1-x)=\frac{2y(1-x-y)}{1-x},
\label{eq:betau}
\end{equation}
and
\begin{align}
\rho_u(L_\perp;u,x,y)
={}&
\mu_k^2(\lambda)
-\lambda(1-\lambda)\left[(x-1+u)(x-1-u)p^2+L_\perp^2\right]
\nonumber\\
&
-\lambda(1-\lambda x)p^2.
\label{eq:rhou}
\end{align}
Consequently,
\begin{equation}
\mathcal I(x,y)=
\Gamma(\epsilon)\,
\frac{\Theta(0\leq y\leq 1-x)}{1-x}
\int_0^1 \dd u
\int \frac{\dd\alpha\,\dd^{d-2}L_\perp}{\ii\pi^{d/2}}
\frac{[\beta_u\alpha+\rho_u-\ii0]^{-\epsilon}}
{[2a_u\alpha+\Lambda_u+\ii0]^2}.
\label{eq:Ialpha}
\end{equation}



\subsubsection{Finite-$\epsilon$ evaluation of the $\alpha$ integral}

Define
\begin{equation}
D_u(\alpha)=2a_u\alpha+\Lambda_u+\ii0,
\qquad
C_u(\alpha)=\beta_u\alpha+\rho_u-\ii0.
\end{equation}
We evaluate
\begin{equation}
J_\eps(u,L_\perp)
\equiv
\Gamma(\eps)
\int_{-\infty}^{\infty}\dd\alpha\,
\frac{C_u(\alpha)^{-\eps}}{D_u(\alpha)^2}.
\label{eq:J-eps-def}
\end{equation}
The two $\alpha$-dependent factors are combined as
\begin{equation}
\frac{\Gamma(\eps)}{D_u^2 C_u^\eps}
=
\Gamma(2+\eps)
\int_0^1\dd r\,
 r(1-r)^{\eps-1}
\frac1{[rD_u+(1-r)C_u]^{2+\eps}}.
\label{eq:alpha-combine}
\end{equation}
The denominator in~\eqref{eq:alpha-combine} is affine in $\alpha$:
\begin{equation}
 rD_u+(1-r)C_u
 =P(r,u)\alpha+Q(r,u,L_\perp)+\ii0_r,
\label{eq:r-denom}
\end{equation}
where
\begin{equation}
 P(r,u)=2r a_u+(1-r)\beta_u,
\label{eq:P-ru}
\end{equation}
and
\begin{equation}
 Q(r,u,L_\perp)=r\Lambda_u+(1-r)\rho_u.
\label{eq:Q-ru}
\end{equation}

\paragraph{Distributional $\alpha$ integral}
Let
\begin{equation}
\mathcal A_\nu(P,Q)
\equiv
\int_{-\infty}^{\infty}\dd\alpha\,
[P\alpha+Q+\ii0]^{-\nu}
\end{equation}
For $P\ne0$, the change of variables $t=P\alpha+Q$ gives a total derivative:
\begin{equation}
\int_{-\infty}^{\infty}\dd t\,(t+\ii0)^{-\nu}
=
-\frac{1}{\nu-1}
\left[(t+\ii0)^{1-\nu}\right]_{-\infty}^{+\infty}=0.
\end{equation}
Thus the $\alpha$ integral vanishes for $P\ne0$.  For $P=0$, it is not an ordinary function but a distribution supported at $P=0$.
To evaluate the $\alpha$ integral we use the Schwinger representation. For $\operatorname{Re}\nu>0$,
\begin{equation}
(x+\ii0)^{-\nu}
=
\frac{e^{-\ii\pi\nu/2}}{\Gamma(\nu)}
\int_0^\infty \dd s\,
 s^{\nu-1}e^{\ii s(x+\ii0)}.
\label{eq:schwinger-i0}
\end{equation}
This follows from the Euclidean Schwinger formula
\begin{equation}
A^{-\nu}=
\frac{1}{\Gamma(\nu)}
\int_0^\infty \dd s\,
 s^{\nu-1}e^{-sA},
\qquad
\operatorname{Re}A>0,
\end{equation}
by setting
\begin{equation}
A=-\ii(x+\ii0).
\end{equation}
Moreover,
\begin{equation}
e^{-sA}=e^{\ii s(x+\ii0)},
\end{equation}
and the phase factor $e^{-\ii\pi\nu/2}$ converts $[-\ii(x+\ii0)]^{-\nu}$ back to $(x+\ii0)^{-\nu}$.
In our case,
\begin{equation}
x=P\alpha+Q.
\end{equation}
Therefore Eq.~\eqref{eq:schwinger-i0} gives
\begin{align}
(P\alpha+Q+\ii0)^{-\nu}
&=
\frac{e^{-\ii\pi\nu/2}}{\Gamma(\nu)}
\int_0^\infty \dd s\,
 s^{\nu-1}e^{\ii s(Q+\ii0)}e^{\ii sP\alpha}\\
 &=
2\pi\delta(P)
\frac{e^{-\ii\pi\nu/2}}{\Gamma(\nu)}
\int_0^\infty \dd s\,
 s^{\nu-2}e^{\ii s(Q+\ii0)}
\end{align}
where we used that the $\alpha$ dependence is now purely Fourier-like
\begin{equation}
\int_{-\infty}^{\infty}\dd\alpha\,e^{\ii sP\alpha}=2\pi\delta(sP).
\label{eq:alpha-fourier-delta}
\end{equation}
and that
\begin{equation}
\delta(sP)=\frac{1}{s}\delta(P),
\qquad s>0.
\end{equation}
The remaining $s$ integral has one lower power of $s$.  Applying the same Schwinger identity with $\nu\to\nu-1$ gives
\begin{equation}
\int_0^\infty \dd s\,
 s^{\nu-2}e^{\ii s(Q+\ii0)}
=
\Gamma(\nu-1)e^{\ii\pi(\nu-1)/2}(Q+\ii0)^{-(\nu-1)}.
\end{equation}
Therefore
\begin{equation}
\mathcal A_\nu(P,Q)
=
2\pi
\frac{\Gamma(\nu-1)}{\Gamma(\nu)}
 e^{-\ii\pi\nu/2}e^{\ii\pi(\nu-1)/2}
\delta(P)(Q+\ii0)^{-(\nu-1)}.
\end{equation}
Since
\begin{equation}
\frac{\Gamma(\nu-1)}{\Gamma(\nu)}=\frac{1}{\nu-1},
\qquad
 e^{-\ii\pi\nu/2}e^{\ii\pi(\nu-1)/2}=e^{-\ii\pi/2}=-\ii,
\end{equation}
we obtain
\begin{equation}
\boxed{
\mathcal A_\nu(P,Q)
=
-\frac{2\pi\ii}{\nu-1}
\delta(P)
(Q+\ii0)^{-(\nu-1)}.
}
\label{eq:A-nu-final}
\end{equation}
For the present integral, $\nu=2+\eps$, so
\begin{equation}
\boxed{
\int_{-\infty}^{\infty}\dd\alpha\,
[P\alpha+Q+\ii0]^{-2-\eps}
=
-\frac{2\pi\ii}{1+\eps}
\delta(P)
(Q+\ii0)^{-1-\eps}.
}
\label{eq:alpha-distribution-final}
\end{equation}



\paragraph{Collapse of the $r$ integral and induced $u$ support}

Applying Eq.~\eqref{eq:alpha-distribution-final} to Eq.~\eqref{eq:alpha-combine} gives
\begin{align}
J_\eps(u,L_\perp)
&=
-\frac{2\pi\ii}{1+\eps}\Gamma(2+\eps)
\int_0^1\dd r\,
 r(1-r)^{\eps-1}
\delta\!\bigl(P(r,u)\bigr)
\bigl[Q(r,u,L_\perp)+\ii0_r\bigr]^{-1-\eps}.
\label{eq:J-after-alpha}
\end{align}
Since
\begin{equation}
P(r,u)=\beta_u+r(2a_u-\beta_u),
\end{equation}
the zero of $P$ is
\begin{equation}
r_\star(u)=\frac{\beta_u}{\beta_u-2a_u}
=\frac{\beta_u}{\beta_u+2(X-u)}.
\label{eq:r-star}
\end{equation}
Moreover,
\begin{equation}
\partial_rP(r,u)=2a_u-\beta_u.
\end{equation}
Therefore
\begin{equation}
\int_0^1\dd r\,f(r)\delta(P(r,u))
=
\begin{cases}
\displaystyle \frac{f(r_\star)}{|\beta_u-2a_u|}, & 0<r_\star<1,\\[1.2em]
0, & r_\star\notin(0,1).
\end{cases}
\label{eq:r-delta-collapse}
\end{equation}
The Feynman parameter $r$ is integrated only over $0\le r\le1$, so the delta function contributes only when its zero lies in this interval.  On the support $0\le y\le1-x$ one has $\beta_u\ge0$.  In the interior, $\beta_u>0$, and the condition $0<r_\star<1$ is equivalent to
\begin{equation}
a_u<0.
\end{equation}
Since $a_u=u-X$, this gives
\begin{equation}
\boxed{0<u<X=1-x.}
\label{eq:u-support-corrected}
\end{equation}
Thus the $r$-integration range induces the remaining support in $u$.

\subsubsection{Current integral after the $\alpha$ and $r$ integrations}

Evaluating the integrand at $r_\star$ gives

\begin{align}
\frac{r_\star(1-r_\star)^{\eps-1}}{|\partial_rP(r_\star,u)|}
&=
\frac{
\dfrac{\beta_u}{\beta_u+2(X-u)}
\left[
\dfrac{2(X-u)}{\beta_u+2(X-u)}
\right]^{\eps-1}
}{\beta_u+2(X-u)}
\nonumber\\
&=
\frac{
\beta_u[2(X-u)]^{\eps-1}
}{[\beta_u+2(X-u)]^{1+\eps}}.
\label{eq:r-weight-jacobian-simplified}
\end{align}
Furthermore,
\begin{align}
Q(r_\star,u,L_\perp)
&=r_\star\Lambda_u+(1-r_\star)\rho_u
\nonumber\\
&=
\frac{\beta_u\Lambda_u+2(X-u)\rho_u}{\beta_u+2(X-u)}
\nonumber\\
&=
\frac{\beta_u\Lambda_u-2a_u\rho_u}{\beta_u+2(X-u)}.
\label{eq:Q-root-before-D}
\end{align}
Define
\begin{equation}
\mathcal D_u(L_\perp)
\equiv
\beta_u\Lambda_u-2a_u\rho_u.
\label{eq:D-u-def}
\end{equation}
Then
\begin{equation}
[Q(r_\star,u,L_\perp)+\ii0]^{-1-\eps}
=
[\beta_u+2(X-u)]^{1+\eps}
[\mathcal D_u(L_\perp)-\ii0]^{-1-\eps}.
\label{eq:Q-cancel-factor}
\end{equation}
 The factor $[\beta_u+2(X-u)]^{1+\eps}$ cancels the denominator in Eq.~\eqref{eq:r-weight-jacobian-simplified}.  Therefore the net result after the $r$ and $\alpha$ integrations is proportional to
\begin{equation}
\beta_u[2(X-u)]^{\eps-1}
[\mathcal D_u(L_\perp)-\ii0]^{-1-\eps}.
\end{equation}
Thus the full integral at this stage is
\begin{equation}
\boxed{
\mathcal I(x,y)
=
\frac{\Theta_{xy}}{X}\,
2\Gamma(1+\eps)
\int_0^X\dd u\,
\beta_u[2(X-u)]^{\eps-1}
\int\frac{\dd^{d-2}L_\perp}{\pi^{(d-2)/2}}
\bigl[\mathcal D_u(L_\perp)-\ii0\bigr]^{-1-\eps}.
}
\label{eq:status-after-alpha-before-Lperp}
\end{equation}
The remaining tasks are the transverse $L_\perp$ integration and the final one-dimensional integration over $u$.

\subsubsection{Transverse integration and one-dimensional representation}

The transverse dependence of $\mathcal D_u$ is linear in $L_\perp^2$.  Using $\beta_u=2X\kappa$, one obtains
\begin{equation}
\mathcal D_u(L_\perp)=2u\kappa L_\perp^2+C_u,
\label{eq:D-u-linear}
\end{equation}
where
\begin{align}
C_u={}&
-2X\kappa\,\mu_l^2(u)-2a_u\,\mu_k^2(\lambda)
\nonumber\\
&+2\Biggl\{
X\kappa\bigl[a_u^2+u(1-u)\bigr]
+a_u\kappa\bigl[X^2-u^2\bigr]
+a_u\frac{y}{1-x}\left(1-\frac{xy}{1-x}\right)
\Biggr\}p^2.
\label{eq:C-u-def}
\end{align}

\textcolor{red}{compare to }
\begin{align}
C_u={}&
-2X\kappa\,\mu_l^2(u)-2a_u\,\mu_k^2(\lambda)
\nonumber\\
&+2\Biggl\{
X\kappa\bigl[a_u^2+u(1-u)\bigr]
+a_u\kappa X^2
+a_u\frac{y}{1-x}\left(1-\frac{xy}{1-x}\right)
\Biggr\}p^2.
\label{eq:C-u-def}
\end{align}

\textcolor{red}{end compare}

The transverse integral yields with $n=d-2=2-2\eps$
\begin{equation}
\int\frac{\dd^n L_\perp}{\pi^{n/2}}
\bigl[2u\kappa L_\perp^2+C_u-\ii0\bigr]^{-1-\eps}
=
\frac{\Gamma(2\eps)}{\Gamma(1+\eps)}
(2u\kappa)^{-1+\eps}(C_u-\ii0)^{-2\eps}.
\label{eq:Lperp-result}
\end{equation}
Substituting Eq.~\eqref{eq:Lperp-result} into Eq.~\eqref{eq:status-after-alpha-before-Lperp}, and using $\beta_u=2X\kappa$, gives
\begin{equation}
\boxed{
\mathcal I(x,y)
=
\Theta_{xy}\,\Gamma(2\eps)(4\kappa)^\eps
\int_0^X\dd u\,
 u^{-1+\eps}(X-u)^{-1+\eps}
(C_u-\ii0)^{-2\eps}.
}
\label{eq:method1-u-representation}
\end{equation}


The endpoint values of $C_u$ are
\begin{equation}
C_0=2X\Delta_1,
\qquad
C_X=-2X\kappa\Delta_0.
\label{eq:C-endpoints}
\end{equation}


\subsubsection{Exact matching of the finite-$\epsilon$ integral}

The one-dimensional integral~\eqref{eq:method1-u-representation} is now mapped to the Euler form from Method 1.  Use
\begin{equation}
 w=\frac{X-u}{X-u+X\kappa},
\label{eq:u-to-w-map}
\end{equation}
or equivalently
\begin{equation}
 u=X-\frac{X\kappa w}{1-w}.
\label{eq:w-to-u-map}
\end{equation}
The endpoints map as
\begin{equation}
 u=X \Longleftrightarrow w=0,
 \qquad
 u=0 \Longleftrightarrow w=\frac1{1+\kappa}.
\label{eq:u-w-endpoints}
\end{equation}
Furthermore,
\begin{equation}
U=1-(1+\kappa)w.
\label{eq:U-under-u-w-map}
\end{equation}
Under this change of variables, Eq.~\eqref{eq:method1-u-representation} becomes
\begin{equation}
\boxed{
\mathcal I_{\rm M2}(x,y)
=
\Gamma(2\eps)
\frac{\Theta_{xy}}{1-x}
\int_0^{1/(1+\kappa)}\dd w\,
 w^{\eps-1}U^{-1+\eps}
[\Delta_x(w)-\ii0]^{-2\eps}.
}
\label{eq:M1-to-M2-w}
\end{equation}
This is precisely the Method-1 representation Eq.~\eqref{eq:method1-w-integral-before-euler}.  Therefore
\begin{equation}
\boxed{
\mathcal I_{\rm M2}(x,y;\eps)
=
\mathcal I_{\rm M1}(x,y;\eps)
}
\label{eq:exact-match}
\end{equation}
as finite-$\eps$ integrals.  Hence this method also yields the exact compact result
\begin{equation}
\boxed{
\mathcal I_{\rm M2}(x,y)
=
\frac{\Theta_{xy}}{1-x}
\Gamma(\eps)^2
(\Delta_0-\ii0)^{-\eps}
(\Delta_1-\ii0)^{-\eps}.
}
\label{eq:M2-final-exact}
\end{equation}

\newpage

\subsection{Use as universal two-loop collinear block}
For a generic hard subgraph $G(x,y)$,
\begin{equation}
I_G^{(2)}=\int_0^1dx\int_0^{1-x}dy\,\mathcal I(x,y)G(x,y).
\label{eq:two-loop-collinear-convolution}
\end{equation}
This is the direct two-loop analogue of the one-loop convolution $I_G=\int dx\,I_c(x)G(x)$.

\newpage
















\end{document}